% Generated draft: requires mathematical and source review.
\documentclass[11pt]{article}
\usepackage[margin=25mm]{geometry}
\usepackage{fontspec}
\setmainfont{Latin Modern Roman}
\usepackage{amsmath,amssymb,mathtools,unicode-math}
\setmathfont{Latin Modern Math}
\usepackage{enumitem,needspace,xurl,hyperref}
\setlength{\parindent}{0pt}
\setlength{\parskip}{5pt}
\setlength{\emergencystretch}{3em}
\newcommand{\N}{\mathbb{N}}
\newcommand{\Z}{\mathbb{Z}}
\newcommand{\Q}{\mathbb{Q}}
\newcommand{\R}{\mathbb{R}}
\newcommand{\C}{\mathbb{C}}
\newcommand{\F}{\mathbb{F}}
\newcommand{\A}{\mathbb{A}}
\newcommand{\PP}{\mathbb{P}}
\newcommand{\E}{\mathbb{E}}
\newcommand{\eps}{\varepsilon}
\newcommand{\dd}{\,\mathrm{d}}
\begin{document}
\section*{UnsolvedMath ID: 700564}
\textbf{Source problem code:} catalogue entry
\subsection{Definitions and mathematical statement}
Let \(k\) be a field with \(\operatorname{char}(k)=0\), and let \(G\) be a \(k\)-powered group. Write \(\mathrm{Rep}_k(G)\) for the family of all canonical matrix representations of \(G\) over \(k\), with both ``\(k\)-powered group'' and ``canonical matrix representation'' understood in the sense of the cited work of Yu. I. Merzlyakov. The supplied record does not give definitions of these terms.

According to the source's account of Merzlyakov's work, \(\mathrm{Rep}_k(G)\) may be regarded as an affine \(k\)-variety. The problem is to give an explicit system of polynomial equations defining this variety, for every field \(k\) of characteristic zero and every \(k\)-powered group \(G\) within the scope of that assertion.

More concretely, if the asserted realization uses an ambient affine space \(\mathbb{A}_k^N\), the requested description would specify its coordinates and an explicit family \(\mathcal E_G\subseteq k[x_1,\ldots,x_N]\) whose common zero locus is the asserted realization of \(\mathrm{Rep}_k(G)\). This coordinate description is conditional: the supplied record specifies neither such an ambient space nor the representation dimensions, and it does not explain whether additional hypotheses on \(G\) or a different convention for affine varieties are used. These qualifications must be resolved from the cited reference before a fully precise coordinate formulation can be given.
\subsection{Short English statement}
For a field \(k\) of characteristic zero and a \(k\)-powered group \(G\), find explicit equations for the affine \(k\)-variety of canonical matrix representations described by Merzlyakov. The definitions of the representation family, its coordinates, and the scope of the asserted variety structure require clarification from the cited reference.
\subsection{Sources}
\begin{itemize}
\item Local UnsolvedMath registry, permanent ID \(700564\), exact source problem code \texttt{catalogue entry}, title ``Kourovka Notebook Problem 9.31.'' The supplied registry labels the problem open; this label is not an independently verified mathematical or literature assessment.
\item \emph{The Kourovka Notebook (Unsolved Problems in Group Theory)}, original problem 9.31, ninth issue (1984); attributed to Yu. I. Merzlyakov. The supplied background also identifies the twenty-first issue edition (2026), \texttt{arXiv:1401.0300v48}. The supplied set description identifies E. I. Khukhro and V. D. Mazurov as editors.
\item Yu. I. Merzlyakov, \emph{Proc. Steklov Inst. Math.} \textbf{167} (1986), 263--266. The source attributes the affine-variety assertion to this reference. No article title or external link is supplied, and the reference has not been independently consulted here.
\end{itemize}
\end{document}
